\documentclass[11pt]{article}
\usepackage[margin=1.1in]{geometry}
\usepackage{amsmath,amssymb,amsthm}
\usepackage{booktabs}
\usepackage[round]{natbib}
\usepackage{hyperref}
\usepackage{microtype}

\newtheorem{proposition}{Proposition}
\newtheorem{remark}{Remark}
\newtheorem{conjecture}{Conjecture}

\title{Dyadic regression}
\author{Peter Cotton}
\date{7 October 2026}

\begin{document}
\maketitle

\begin{abstract}
Dyadic regression forecasts a series from its last value, the sum of the two values before that,
the sum of the four before those, and so on, with one coefficient per block. With $L$ lags it fits
$\log_2(L+1)$ coefficients. It represents every trailing moving average of length $2^m-1$ exactly,
and it approximates any power-law lag profile $k^{-\alpha}$ to within relative error $2^{\alpha/2}-1$
uniformly in $L$. Weighting each lag by $1/(2k-1)$ inside its block cuts that error to about
$2^{|\alpha-1|/2}-1$, from $0.62$ to $0.15$ for fractional differencing at $d=0.4$, and makes the
all-equal point exactly the harmonic forecaster of \citet{fishelson2026}.
\end{abstract}

\section{The model}

Partition the lags $k=1,\dots,L$, with $L=2^J-1$, into dyadic blocks
\[
G_j=\{2^j,\dots,2^{j+1}-1\},\qquad j=0,\dots,J-1,
\]
so $G_0=\{1\}$, $G_1=\{2,3\}$, $G_2=\{4,\dots,7\}$. Dyadic regression is the forecast
\begin{equation}\label{eq:model}
\hat y_{t+1}=\sum_{j=0}^{J-1}\theta_j\,S_j(t),\qquad S_j(t)=\sum_{k\in G_j}y_{t+1-k},
\end{equation}
an autoregression whose coefficient on lag $k$ is $\theta_{j(k)}$, shared within each block. Writing
the regressors as block averages $\bar S_j=S_j/2^j$ instead gives the same model with coefficients
$2^j\theta_j$. The \texttt{grouped\_ar} transform of the \texttt{skaters} library implements
\eqref{eq:model} with block sums, fitted by recursive least squares with exponential forgetting and
a ridge toward zero.

The number of coefficients is $J=\log_2(L+1)$, against $L$ for an unrestricted autoregression. The
question is what the restriction costs and what it buys.

\section{What the blocks can represent}

\paragraph{Moving averages.} The trailing average of length $W=2^m-1$ is exactly
$\theta_j=1/W$ for $j<m$ and $\theta_j=0$ for $j\ge m$. Every other window length differs from one of
these by part of a single block. Dyadic regression therefore contains, exactly or nearly, the moving
average at every time scale, with one coefficient per doubling. This is the same covering that lets
online learners compete with the best window of any length using $O(\log T)$ experts
\citep{daniely2015}.

\paragraph{Power-law profiles.} Long memory shows up as autoregressive coefficients that decay like a
power of the lag. For fractional differencing $(1-B)^d y_t=\varepsilon_t$ with $0<d<\tfrac12$ the
coefficients of the AR($\infty$) representation satisfy $\pi_k\sim C\,k^{-(1+d)}$.

\begin{proposition}\label{prop:powerlaw}
Let $\varphi_k=c\,k^{-\alpha}$ with $\alpha>0$. Choosing $\theta_j$ as the geometric mean of
$\varphi$ at the two ends of block $G_j$ gives, for every $L$ and every lag $k\le L$,
\[
\Bigl|\frac{\theta_{j(k)}}{\varphi_k}-1\Bigr|\;\le\;2^{\alpha/2}-1 .
\]
Consequently, for any series, the one-step forecast of the dyadic model differs from that of the
full AR($L$) model by at most $(2^{\alpha/2}-1)\sum_{k\le L}|\varphi_k|\,|y_{t+1-k}|$.
\end{proposition}

\begin{proof}
On $G_j$ the ratio of the largest to the smallest value of $\varphi$ is
$\bigl((2^{j+1}-1)/2^j\bigr)^{\alpha}<2^{\alpha}$. The geometric mean $\theta_j$ of the two extremes
is within a factor $2^{\alpha/2}$ of every value in the block, so
$\theta_j/\varphi_k\in[2^{-\alpha/2},2^{\alpha/2}]$, and $1-2^{-\alpha/2}\le 2^{\alpha/2}-1$. The
forecast bound follows from the triangle inequality.
\end{proof}

The bound does not shrink as $L$ grows, and it is not small: $0.23$ at $\alpha=0.6$, $0.41$ at
$\alpha=1$, and $0.62$ for fractional differencing at $d=0.4$. It is a worst case over lags, and the
least-squares coefficients of a fitted model trade it against the variance saving of $L-J$ fewer
parameters. Blocks growing by a ratio $r<2$ instead of $2$ cut the error to $r^{\alpha/2}-1$ at the
cost of $\log_r(L)$ coefficients. Table~\ref{tab:powerlaw} gives the numbers.

\begin{table}[h]
\centering
\begin{tabular}{lccc}
\toprule
profile & $r=2$ (13 blocks) & $r=1.5$ (21) & $r=1.25$ (38)\\
\midrule
$k^{-0.6}$ & 0.231 & 0.129 & 0.069\\
$k^{-1}$ & 0.414 & 0.225 & 0.118\\
$k^{-1.4}$ & 0.624 & 0.328 & 0.169\\
$(1-B)^{0.2}$, $\pi_k$ & 0.516 & & \\
$(1-B)^{0.4}$, $\pi_k$ & 0.624 & & \\
\bottomrule
\end{tabular}
\caption{Worst relative error of the best one-coefficient-per-block approximation over lags
$k\le 4096$, for geometric blocks of ratio $r$. Dyadic regression is $r=2$. The fractional-differencing
rows use the exact AR($\infty$) coefficients and match the bound of
Proposition~\ref{prop:powerlaw} with $\alpha=1+d$. Computed by \texttt{check\_dyadic.py}.}
\label{tab:powerlaw}
\end{table}

This is the dyadic form of the argument for the heterogeneous autoregressive model of
\citet{corsi2009}, which regresses on daily, weekly and monthly averages and reproduces long memory
with three coefficients. \citet{hwang2014} show that the infinite-order version is a long-memory
model exactly when its coefficients decrease exponentially across horizons, and that a fitted
finite-order version loses little to truncation. Dyadic regression fixes the horizons at the powers
of two and lets the data choose the coefficients.

\section{Harmonic dyadic regression}

Most of the error in Proposition~\ref{prop:powerlaw} comes from treating every lag in a block alike
when the profile falls across the block. Weighting lag $k$ by $1/(2k-1)$ inside its block,
\begin{equation}\label{eq:harm}
\hat y_{t+1}=\sum_{j=0}^{J-1}\theta_j\,H_j(t),\qquad H_j(t)=\sum_{k\in G_j}\frac{y_{t+1-k}}{2k-1},
\end{equation}
keeps one coefficient per block and puts the profile's typical decay into the regressor. A power
law $c\,k^{-\alpha}$ then has to be matched only up to the residual factor $k^{1-\alpha}$, which
varies by at most $2^{|\alpha-1|}$ across a block. The argument of Proposition~\ref{prop:powerlaw}
gives a worst relative error of $2^{|\alpha-1|/2}-1$ for weights $1/k$. The weights $1/(2k-1)$ add a
small correction at short lags, $0.054$ at $\alpha=1$ where the bound is zero. For long memory,
$\alpha=1+d$, the error is $2^{d/2}-1$: $0.035$ at $d=0.1$ and $0.149$ at $d=0.4$, against $0.46$
and $0.62$ for plain dyadic regression. Table~\ref{tab:harm} gives the numbers.

\begin{table}[h]
\centering
\begin{tabular}{lcc}
\toprule
profile & dyadic & harmonic dyadic\\
\midrule
$k^{-0.6}$ & 0.231 & 0.152\\
$k^{-1}$ & 0.414 & 0.054\\
$k^{-1.4}$ & 0.624 & 0.149\\
$k^{-2}$ & 1.000 & 0.414\\
$(1-B)^{0.1}$ & 0.464 & 0.035\\
$(1-B)^{0.2}$ & 0.516 & 0.072\\
$(1-B)^{0.3}$ & 0.569 & 0.109\\
$(1-B)^{0.4}$ & 0.624 & 0.149\\
\bottomrule
\end{tabular}
\caption{Worst relative error over lags $k\le 4096$, thirteen blocks in both columns. Computed by
\texttt{check\_dyadic.py}.}
\label{tab:harm}
\end{table}

With every $\theta_j$ equal, \eqref{eq:harm} is a weighted average of past values with weight
$1/(2k-1)$ on lag $k$. That is exactly the stored forecast of \citet{fishelson2026}, discussed next.

\section{The equal-weight point}

Set every coefficient on the block averages equal, which is $\theta_j\propto 2^{-j}$ on the sums.
Lag $k$ in block $j$ then gets weight $w_k=2^{-j}$, and since $2^j\le k<2^{j+1}$,
\[
\frac1k\;\le\;w_k\;<\;\frac2k .
\]
The profile is $1/k$ up to a factor of two, a staircase with one step per doubling. Every dyadic
scale gets the same total weight, and the weights over $L$ lags sum to $\log_2(L+1)$.

\citet{fishelson2026} forecast with exactly this shape. Their algorithm stores a weighted average of
past outcomes with weights $1/(2k-1)$ by age, and announces a mixture of its past stored averages
with the same weights. They prove that it is calibrated against every outcome sequence, with
average calibration error falling like $1/\log T$, and that the kernel $1/(2(t-s)-1)$, the discrete
Hilbert transform, is the best choice up to a constant. The reason is that a forecast's calibration
error is a comparison between the outcomes that built it and the outcomes on which it is used, and
a kernel with equal weight at every scale offers an adversary no scale to attack.

For harmonic dyadic regression the match is exact. With all $\theta_j$ equal and normalized,
\eqref{eq:harm} is their stored average $p_t$. Their calibrated forecast is one step further, the
same harmonic mixture of the stored averages $p_s$, $s<t$, and the guarantee is for that mixture.
For plain dyadic regression the equal-weight point is the staircase above.

\begin{conjecture}
The \citet{fishelson2026} guarantee holds, with a larger constant, when their harmonic kernel is
replaced by the dyadic staircase.
\end{conjecture}

Their own remark that the weights $1/k$ give the same rates as $1/(2k-1)$ suggests the result does
not depend on the exact kernel, and dyadic analogues of the Hilbert transform are classically
bounded. Their proof, however, uses moment bounds specific to the Hilbert matrix, so the staircase
needs its own check. Harmonic dyadic regression avoids the question.

\begin{remark}[Fitted versus frozen]
Fitting the coefficients is what makes dyadic regression forecast well on real series, where some
scales predict and others do not. Freezing them at the equal-weight point, and announcing the harmonic mixture of
the stored averages rather than any single average, is what makes the \citet{fishelson2026}
forecaster robust, since a learned preferred scale is exactly what an adversary exploits. Averaging
that mixture into one forecast loses the guarantee. A frozen forecaster that uses only past outcomes is calibrated but carries no
information about the input. The two can be combined by shrinking the fitted coefficients toward
the equal-weight point instead of toward zero, so the model starts from the forecaster no sequence
can exploit and moves only as far as the data justify. Whether the shrunk model keeps any form of the
guarantee is open.
\end{remark}

\section{Summary}

Dyadic regression is an autoregression with one coefficient per doubling of the lag. It contains
the moving average at every scale and approximates power-law memory to a fixed relative error with
$\log_2 L$ coefficients. Weighting lags by $1/(2k-1)$ inside each block keeps the same number of
coefficients, cuts the long-memory error by a factor of four or more, and makes the all-equal point
exactly the scale-free forecaster that online calibration theory singles out. The numbers in this note are reproduced by \texttt{check\_dyadic.py}.

\begin{thebibliography}{5}
\bibitem[Corsi(2009)]{corsi2009} Corsi, F. (2009). A simple approximate long-memory model of
  realized volatility. \emph{Journal of Financial Econometrics} 7(2), 174--196.
\bibitem[Daniely et al.(2015)]{daniely2015} Daniely, A., Gonen, A. and Shalev-Shwartz, S. (2015).
  Strongly adaptive online learning. \emph{Proceedings of the 32nd International Conference on
  Machine Learning}, PMLR 37.
\bibitem[Fishelson and Mohri(2026)]{fishelson2026} Fishelson, M. and Mohri, M. (2026).
  High-dimensional online calibration from harmonic weights. arXiv:2610.07740.
\bibitem[Hwang and Shin(2014)]{hwang2014} Hwang, E. and Shin, D.~W. (2014). Infinite-order,
  long-memory heterogeneous autoregressive models. \emph{Computational Statistics and Data
  Analysis} 76, 339--358.
\end{thebibliography}

\end{document}
